
\subsubsection{2.1 FAIR Data Principles and Keyword Findability}

The FAIR Guiding Principles \citep{Wilkinson2016FAIR} established a framework for scientific data management organized around Findability (F), Accessibility (A), Interoperability (I), and Reusability (R). The F1 sub-principle --- that data should be assigned globally unique persistent identifiers and described with rich machine-actionable metadata --- is directly operationalizable in ML repositories via keyword tagging. The Wilkinson et al. framework has accumulated over 15,976 citations and has influenced repository design and curation policy broadly.

The FAIR literature has, however, remained primarily normative. Papers describe desirable metadata properties and propose compliance checklists without measuring the quantitative effect of specific FAIR F1 interventions on dataset adoption outcomes. The Croissant-RAI specification \citep{Jain2024Croissant} --- developed partly by OpenML's founder --- proposes machine-readable structured metadata for ML datasets and identifies keyword tags as a primary findability mechanism, but does not provide empirical IRR estimates. Trišović et al. [2025] study automated FAIR compliance scoring but measure compliance richness rather than adoption outcomes. The present study fills this gap by providing the first negative binomial regression IRR for keyword tag presence on a major ML platform.

\subsubsection{2.2 ML Dataset Adoption and Metadata Quality}

Yang et al. [2024] study the relationship between documentation quality (dataset cards on HuggingFace) and dataset popularity, finding that richer documentation predicts higher download counts. This is a parallel finding --- structured metadata presence predicts platform adoption --- but differs from the present study in three important respects: the platform (HuggingFace, a model-centric platform with different search mechanics), the metadata type (free-form dataset cards rather than keyword tags), and the adoption metric (download counts rather than ML task registrations). Yang et al. do not isolate keyword tagging as a specific FAIR F1 operand, nor do they apply negative binomial regression to handle overdispersion.

Lachmuth et al. [2025] study FAIR metadata compliance and dataset reuse in the BonaRes agricultural data repository, finding that compliance predicts reuse (measured as paper citations). This provides cross-domain corroboration of the FAIR-adoption linkage, but in a domain-specific repository with different platform mechanics and a different adoption proxy. The tag-indexed search mechanism characterizing OpenML does not directly apply to BonaRes's architecture.

Oreamuno et al. [2023] survey documentation practices across ML datasets and find that poor tagging leads to poor discoverability --- consistent with the present findings --- but do not provide regression estimates. Chapman et al. [2019] survey dataset search behavior and identify keywords as the primary mechanism by which researchers discover datasets, providing theoretical grounding for the FAIR F1 mechanism hypothesis. Afzal et al. [2020] propose a data readiness framework incorporating metadata richness but focus on quality assessment rather than adoption regression.

Orr and Crawford [2024] situate ML dataset metadata in sociological context, arguing that metadata reflects curatorial effort and community norms rather than purely technical properties --- a perspective consistent with the present finding that tagging behavior is bimodal on OpenML.

No prior study provides an NB-2 incidence rate ratio for keyword tag presence specifically, predicting ML task adoption, with overdispersion testing, on an ML dataset repository.

\subsubsection{2.3 Negative Binomial Regression for Platform Adoption}

Count data regression for platform adoption outcomes requires careful model selection. Standard Poisson regression assumes mean-variance equality; real-world count outcomes such as task registrations exhibit overdispersion (variance substantially exceeding mean). The Negative Binomial Type 2 (NB-2) model accommodates overdispersion via a quadratic variance function [Cameron and Trivedi, 1986; 2013].

The Cameron-Trivedi Lagrange multiplier test (CT LR) provides a formal overdispersion test. In the present setting, CT LR=7,356.36 far exceeds the critical value ($\chi^{2}(1)=3.84)$, confirming NB-2 as the appropriate model family. Cabansag and Ntegeka [2026] demonstrate the NB-2 IRR interpretation in a related count regression context, providing methodological precedent.

